% Template for post or presentation abstracts submitted to ILAS 2022
% 1. Fill in the presenter's information (name, affiliation, email
% address) and talk information (title of presentation, abstract).
% 2. Compile to make sure there are no errors.
% 3. Save the .tex file for your abstract with a distinctive name,
% ideally "FamilynameGivenname.tex".
% 4. Upload your .tex file to https://clr.ie/129557
%
% * Please keep your LaTeX commands simple, and avoid the use of
%    user-defined macros.
% * Include details of co-authors and funding acknowledgements after the abstract.
% * Upload only the LaTeX file (with .tex extension).
% * Questions: Contact Niall Madden (Niall.Madden@NUIGalway.ie)

\documentclass[a4paper]{amsart}
\usepackage{amssymb}
\usepackage{hyperref}
\pagestyle{empty}
\date{}

%% IGNORE THE NEXT 4 LINES
\makeatletter %
\def\labelenumi{\theenumi.}
\def\theenumi{\@arabic\c@enumi}
\makeatother
%% STOP IGNORING NOW

\begin{document}

\title{Linearizations of matrix polynomials via Rosenbrock polynomial system matrices}
\author{Froil\'an Dopico} %Presenting author only
\address{Universidad Carlos III de Madrid, Spain} % Affiliation of presenting author
\email{dopico@math.uc3m.es} % Email address of presenting author only

\maketitle

% The abstract  ***no more than one page***

In the seventies, Rosenbrock \cite{Ros70} introduced the concept of a polynomial system matrix $L(\lambda)$ of an arbitrary rational matrix $R(\lambda)\in\mathbb{F}(\lambda)^{m\times n}$, where $\mathbb{F}$ is an arbitrary field.
Such a system matrix is partitioned in a quadruple $\{A(\lambda),B(\lambda),\allowbreak C(\lambda),D(\lambda)\}$ of compatible matrix polynomials 
\[ L(\lambda):=\left[\begin{array}{cr} A(\lambda) & -B(\lambda) \\ C(\lambda) & D(\lambda)
\end{array}\right]
\]
such that its Schur complement with respect to $D(\lambda)$ equals $R(\lambda)$. That is, $R(\lambda)= D(\lambda) + C(\lambda)A(\lambda)^{-1}B(\lambda)$. Since then, the concept of polynomial system matrix has played a key role in linear system theory and control theory. Later, in a fully independent way, I. Gohberg, M. A. Kaashoek, P. Lancaster, and L. Rodman introduced the concepts of linearization \cite{GohbergLancasterRodman09} and strong linearization \cite{gkl} of matrix polynomials. The concepts of linearization and strong linearization of matrix polynomials have been widely used in the last two decades, both from theoretical and numerical perspectives, by many authors all over the world, and many explicitly constructible classes of linearizations and strong linearizations have been developed based on them. In this talk, we prove that some of the most important classes of linearizations of matrix polynomials are, modulo block permutations, linear polynomial system matrices whose matrix $A(\lambda)$ is unimodular.

%% Coauthor details here (optional - delete if not applicable)
% and funding acknowledgment (optional - delete if not applicable)
\emph{This is joint work with Silvia Marcaida (Universidad del Pa\'is Vasco
UPV/EHU, Spain), Mar\'ia del Carmen Quintana (Aalto University, Finland) and Paul Van Dooren (Université catholique de Louvain, Belgium). \\
This work is part of the ``Proyecto de I+D+i PID2019-106362GB-I00 financiado
por  MCIN/AEI/10.13039/501100011033''.}

\begin{thebibliography}{1}

\bibitem{gkl}
I. Gohberg, M. A. Kaashoek, and P. Lancaster.
\newblock General theory of regular matrix polynomials and band Toeplitz operators.
\newblock {\em Integral Equations Operator Theory} 11:776-882, (1988).

\bibitem{GohbergLancasterRodman09} I. Gohberg, P. Lancaster, and L. Rodman.
\newblock Matrix Polynomials.
\newblock SIAM Publications, 2009.
\newblock Originally published: Academic Press, New York, 1982.

\bibitem{Ros70} H.H. Rosenbrock.
\newblock State-Space and Multivariable Theory.
\newblock Thomas Nelson and Sons, London, 1970.
\end{thebibliography}


\end{document}


% Please leave the next bit alone
%%% Local Variables:
%%% mode: latex
%%% TeX-master: "../ILAS-2022"
%%% End:
